\textbf{Phần 1. Trắc nghiệm lựa chọn}
\setcounterpageref{ex}{0}

\begin{ex} %[2D1V04]
Năm nay Mai $8$ tuổi, mẹ Mai $34$ tuổi. Hỏi sau $5$ năm nữa mẹ hơn Mai bao nhiêu tuổi?
\choice
{$31$ tuổi}
{$21$ tuổi}
{$39$ tuổi}
{\True $26$ tuổi}
\loigiai{
Hiệu số tuổi của hai người không bao giờ thay đổi theo thời gian. Mẹ hơn Mai: $34 - 8 = 26$ (tuổi).
\textbf{Đáp án của câu hỏi là D}
}
\end{ex}

\begin{ex} %[2D1V04]
Lớp 2A có $35$ học sinh, lớp 2B có ít hơn lớp 2A là $4$ học sinh. Cả hai lớp có tất cả bao nhiêu học sinh?
\choice
{$62$ học sinh}
{\True $66$ học sinh}
{$70$ học sinh}
{$68$ học sinh}
\loigiai{
Số học sinh lớp 2B là: $35 - 4 = 31$ (học sinh).\nCả hai lớp có: $35 + 31 = 66$ (học sinh).
\textbf{Đáp án của câu hỏi là B}
}
\end{ex}

\begin{ex} %[2D1V04]
Đoạn thẳng $AB$ dài $1\text{ dm } 4\text{ cm}$, đoạn thẳng $CD$ dài $9\text{ cm}$. Đoạn thẳng $AB$ dài hơn đoạn thẳng $CD$ bao nhiêu xăng-ti-mét?
\choice
{$7\text{ cm}$}
{$4\text{ cm}$}
{\True $5\text{ cm}$}
{$6\text{ cm}$}
\loigiai{
Đổi $1\text{ dm } 4\text{ cm} = 14\text{ cm}$.\nĐoạn thẳng $AB$ dài hơn đoạn thẳng $CD$ là: $14 - 9 = 5\text{ (cm)}$.
\textbf{Đáp án của câu hỏi là C}
}
\end{ex}

\begin{ex} %[2D1V04]
Một đàn gà có $45$ con gà mái. Số gà trống ít hơn số gà mái $23$ con. Hỏi đàn gà có tất cả bao nhiêu con gà?
\choice
{\True $67$ con gà}
{$68$ con gà}
{$22$ con gà}
{$65$ con gà}
\loigiai{
Số gà trống là: $45 - 23 = 22$ (con).\nCả đàn gà có: $45 + 22 = 67$ (con).
\textbf{Đáp án của câu hỏi là A}
}
\end{ex}

\begin{ex} %[2D1V04]
Trong một hàng dọc, Nam đứng thứ $12$ tính từ đầu hàng và đứng thứ $8$ tính từ cuối hàng. Hỏi hàng đó có tất cả bao nhiêu bạn?
\choice
{$20$ bạn}
{$18$ bạn}
{$21$ bạn}
{\True $19$ bạn}
\loigiai{
Số bạn trong hàng là: $12 + 8 - 1 = 19$ (bạn).
\textbf{Đáp án của câu hỏi là D}
}
\end{ex}

\begin{ex} %[2D1V04]
Thùng thứ nhất có $42\text{ L}$ dầu, thùng thứ hai nhiều hơn thùng thứ nhất $15\text{ L}$ dầu. Cả hai thùng có tất cả bao nhiêu lít dầu?
\choice
{$89\text{ L}$}
{\True $99\text{ L}$}
{$97\text{ L}$}
{$100\text{ L}$}
\loigiai{
Thùng thứ hai có: $42 + 15 = 57\text{ (L)}$.\nCả hai thùng có: $42 + 57 = 99\text{ (L)}$.
\textbf{Đáp án của câu hỏi là B}
}
\end{ex}

\begin{ex} %[2D1V04]
Bao ngô nặng $48\text{ kg}$, bao thóc nhẹ hơn bao ngô $12\text{ kg}$. Hỏi cả hai bao nặng bao nhiêu ki-lô-gam?
\choice
{$72\text{ kg}$}
{$96\text{ kg}$}
{\True $84\text{ kg}$}
{$86\text{ kg}$}
\loigiai{
Bao thóc nặng: $48 - 12 = 36\text{ (kg)}$.\nCả hai bao nặng: $48 + 36 = 84\text{ (kg)}$.
\textbf{Đáp án của câu hỏi là C}
}
\end{ex}

\begin{ex} %[2D1V04]
Có hai ngăn sách, ngăn trên có $38$ quyển sách. Nếu chuyển $5$ quyển từ ngăn trên xuống ngăn dưới thì ngăn trên vẫn nhiều hơn ngăn dưới $4$ quyển. Hỏi ban đầu ngăn trên hơn ngăn dưới bao nhiêu quyển?
\choice
{\True $14$ quyển}
{$9$ quyển}
{$10$ quyển}
{$12$ quyển}
\loigiai{
Chuyển 5 quyển làm hiệu giảm đi $5 \times 2 = 10$ quyển. Lúc này còn hơn 4 quyển nên ban đầu hơn: $10 + 4 = 14$ quyển.
\textbf{Đáp án của câu hỏi là A}
}
\end{ex}

\begin{ex} %[2D1V04]
Một cây cau cao $8\text{ m}$, cây dừa cao hơn cây cau $3\text{ m}$, cây xoài thấp hơn cây dừa $4\text{ m}$. Chiều cao của cây xoài là
\choice
{$6\text{ m}$}
{$8\text{ m}$}
{$9\text{ m}$}
{\True $7\text{ m}$}
\loigiai{
Cây dừa cao: $8 + 3 = 11\text{ (m)}$.\nCây xoài cao: $11 - 4 = 7\text{ (m)}$.
\textbf{Đáp án của câu hỏi là D}
}
\end{ex}

\begin{ex} %[2D1V04]
Lan có nhiều hơn Huệ $16$ bông hoa. Lan tặng Huệ $5$ bông hoa. Hỏi bây giờ Lan còn nhiều hơn Huệ bao nhiêu bông hoa?
\choice
{$11$ bông hoa}
{\True $6$ bông hoa}
{$8$ bông hoa}
{$10$ bông hoa}
\loigiai{
Khi Lan tặng Huệ 5 bông hoa thì số hoa Lan hơn Huệ giảm đi: $5 \times 2 = 10$ (bông).\nBây giờ Lan còn nhiều hơn Huệ: $16 - 10 = 6$ (bông hoa).
\textbf{Đáp án của câu hỏi là B}
}
\end{ex}

\begin{ex} %[2D1V04]
Tìm hai số có hiệu bằng $18$, biết rằng nếu thêm vào số bị trừ $6$ đơn vị thì được số mới bằng $45$.
\choice
{$45\text{ và } 27$}
{$38\text{ và } 20$}
{\True $39\text{ và } 21$}
{$40\text{ và } 22$}
\loigiai{
Số bị trừ là: $45 - 6 = 39$. Số trừ là: $39 - 18 = 21$.
\textbf{Đáp án của câu hỏi là C}
}
\end{ex}

\begin{ex} %[2D1V04]
Một cửa hàng buổi sáng bán được $56$ quả dưa hấu, buổi chiều bán được ít hơn buổi sáng $14$ quả dưa hấu. Cả hai buổi cửa hàng bán được bao nhiêu quả?
\choice
{\True $98$ quả}
{$96$ quả}
{$84$ quả}
{$102$ quả}
\loigiai{
Buổi chiều bán được: $56 - 14 = 42$ (quả).\nCả hai buổi bán được: $56 + 42 = 98$ (quả).
\textbf{Đáp án của câu hỏi là A}
}
\end{ex}

\textbf{Phần 2. Trắc nghiệm đúng sai}
\setcounterpageref{ex}{0}

\begin{ex} %[2D1V04]
Xét tính đúng sai của các bài toán về nhiều hơn, ít hơn:
\choiceTF
{\True Hiệu số tuổi của hai người không đổi theo thời gian.}
{Nếu Nam nhiều hơn Bình 8 viên bi, Nam cho Bình 4 viên bi thì Bình sẽ nhiều hơn Nam.}
{\True Một đoạn dây dài 20 cm bị cắt thành 2 đoạn, đoạn thứ nhất dài hơn đoạn thứ hai 4 cm thì đoạn dài là 12 cm.}
{\True Muốn tìm số lớn ta lấy số bé cộng với phần hơn.}
\loigiai{
\begin{itemize}
    \item Ý a) Đúng: Mỗi năm mỗi người đều tăng thêm 1 tuổi nên hiệu số tuổi luôn giữ nguyên.
    \item Ý b) Sai: Khi Nam cho Bình 4 viên bi thì hai bạn có số bi bằng nhau: $8 - 4 = 4$ và bên kia $0 + 4 = 4$.
    \item Ý c) Đúng: Đoạn dài là: $(20 + 4) : 2 = 12\text{ cm}$.
    \item Ý d) Đúng: Số lớn = Số bé + Phần hơn.
\end{itemize}
\begin{itemize}
    \item a) \textbf{Đúng}.
    \item b) \textbf{Sai}.
    \item c) \textbf{Đúng}.
    \item d) \textbf{Đúng}.
\end{itemize}
}
\end{ex}

\begin{ex} %[2D1V04]
Lớp 2A có 32 học sinh, lớp 2B có 36 học sinh:
\choiceTF
{\True Lớp 2B nhiều hơn lớp 2A là 4 học sinh.}
{Chuyển 4 học sinh từ lớp 2B sang lớp 2A thì hai lớp bằng nhau.}
{\True Cả hai lớp có tất cả 68 học sinh.}
{Nếu lớp 2A có thêm 2 bạn và lớp 2B bớt 2 bạn thì lớp 2B vẫn nhiều hơn.}
\loigiai{
\begin{itemize}
    \item Ý a) Đúng: $36 - 32 = 4$ học sinh.
    \item Ý b) Sai: Chỉ cần chuyển $4 : 2 = 2$ học sinh thì cả hai lớp đều có 34 học sinh.
    \item Ý c) Đúng: $32 + 36 = 68$ học sinh.
    \item Ý d) Sai: Khi đó lớp 2A có 34 bạn, lớp 2B có 34 bạn, hai lớp bằng nhau.
\end{itemize}
\begin{itemize}
    \item a) \textbf{Đúng}.
    \item b) \textbf{Sai}.
    \item c) \textbf{Đúng}.
    \item d) \textbf{Sai}.
\end{itemize}
}
\end{ex}

\textbf{Phần 3. Trả lời ngắn}
\setcounterpageref{ex}{0}

\begin{ex} %[2D1V04]
Bác Ba nuôi 38 con vịt và nuôi số gà nhiều hơn số vịt 15 con. Hỏi bác Ba nuôi tất cả bao nhiêu con gà và vịt?
\shortans{91}
\loigiai{
\textbf{Đáp số:} $91$\\
Số gà: $38 + 15 = 53$ con. Tổng số: $38 + 53 = 91$ con.
}
\end{ex}

\begin{ex} %[2D1V04]
Hiện nay anh 14 tuổi, em 8 tuổi. Hỏi khi em bằng tuổi anh hiện nay thì anh bao nhiêu tuổi?
\shortans{20}
\loigiai{
\textbf{Đáp số:} $20$\\
Anh hơn em: $14 - 8 = 6$ tuổi. Khi em 14 tuổi thì anh: $14 + 6 = 20$ tuổi.
}
\end{ex}

\begin{ex} %[2D1V04]
Túi thứ nhất có 35 viên kẹo, túi thứ hai có ít hơn túi thứ nhất 9 viên kẹo. Hỏi cả hai túi có tất cả bao nhiêu viên kẹo?
\shortans{61}
\loigiai{
\textbf{Đáp số:} $61$\\
Túi 2: $35 - 9 = 26$ viên. Cả hai túi: $35 + 26 = 61$ viên.
}
\end{ex}

\begin{ex} %[2D1V04]
Hùng có nhiều hơn Dũng 18 viên bi. Hùng phải cho Dũng bao nhiêu viên bi để hai bạn có số bi bằng nhau?
\shortans{9}
\loigiai{
\textbf{Đáp số:} $9$\\
$18 : 2 = 9$ viên bi.
}
\end{ex}

\begin{ex} %[2D1V04]
Một mảnh vải dài 45 dm, mảnh vải thứ hai ngắn hơn mảnh vải thứ nhất 18 dm. Hỏi cả hai mảnh vải dài bao nhiêu đề-xi-mét?
\shortans{72}
\loigiai{
\textbf{Đáp số:} $72$\\
Mảnh 2: $45 - 18 = 27\text{ dm}$. Tổng: $45 + 27 = 72\text{ dm}$.
}
\end{ex}

\begin{ex} %[2D1V04]
Trong một cuộc thi chạy, từ vị trí của Mai về đích còn 25 m, Mai chạy trước Lan 8 m. Hỏi từ vị trí của Lan về đích còn bao nhiêu mét?
\shortans{33}
\loigiai{
\textbf{Đáp số:} $33$\\
Vị trí của Lan cách đích: $25 + 8 = 33\text{ m}$.
}
\end{ex}
